\DocumentMetadata{pdfversion=1.7,pdfstandard=A-3a,tagging=on,lang=zh-CN}
\documentclass[12pt,fontset=windows]{ctexart}

\usepackage[colorlinks,allcolors=blue]{hyperref}
\usepackage[a4paper,hmargin=2cm,vmargin=2cm]{geometry}
\usepackage{xparse}
\usepackage{unicode-math}
\usepackage{amsmath}
\usepackage[backend=biber,style=gb7714-2015,autocite=superscript]{biblatex}
\usepackage{fancyhdr}
\usepackage{setspace}

% 字体
% \setmainfont{Times New Roman}
% \setCJKmainfont{NSimSun}[AutoFakeBold=1.5]
\newCJKfontfamily{\xingkai}{STXingkai}[Renderer=HarfBuzz]

% 布局
\pagestyle{fancy}
\fancyhf{}
\renewcommand{\headrulewidth}{0cm}
\fancyfoot[C]{\thepage}

% 封面
\makeatletter
\renewcommand{\maketitle}{
\begin{titlepage}
	\centering

	\vspace{5ex}

	\makebox[15em][s]{\zihao{3} \xingkai 西北农林科技大学}

	\vspace{10ex}

	{\zihao{2} 研究生课程论文}

	\vspace{12ex}

	{\zihao{-4}（课程名称：\underline{《中国特色社会主义理论与实践研究》}）}

	\vfill

	\begin{spacing}{1.16}\zihao{-3}\setlength{\parskip}{8pt}
		\makebox[6em][s]{年级、姓名}：\underline{\makebox[20em][c]{25级 \@author}}

		\makebox[6em][s]{学院、专业}：\underline{\makebox[20em][c]{机械与电子工程学院 机械专业}}

		\makebox[6em][s]{结业成绩}：\underline{\makebox[20em][c]{}}

		\makebox[6em][s]{结业学期}：\underline{\makebox[20em][c]{2025年秋季学期}}

		\makebox[6em][s]{任课老师姓名}：\underline{\makebox[20em][c]{方建斌}}
	\end{spacing}

\end{titlepage}

}
\makeatother

% 标题
\makeatletter
\newcommand{\articleTitle}{\begin{center}\zihao{3}\heiti\@title\end{center}}
\makeatother


% 摘要
\renewenvironment{abstract}{

	\begin{center}
		\Large\bfseries
		摘\hspace{1em}要
	\end{center}

	\normalfont\normalsize

}{}

% 关键词
\NewDocumentCommand{\keywords}{m o o o o o o o o}{

	{
		{\bfseries 关键词} #1\IfNoValueTF{#2}{}{；#2}\IfNoValueTF{#3}{}{；#3}\IfNoValueTF{#4}{}{；#4}\IfNoValueTF{#5}{}{；#5}\IfNoValueTF{#6}{}{；#6}\IfNoValueTF{#7}{}{；#7}\IfNoValueTF{#8}{}{；#8}\IfNoValueTF{#9}{}{；#9}
	}

	\normalfont\normalsize

}
\addbibresource{./bibliography.bib}

\title{新质生产力驱动下“十五五”农业机械化与农业现代化融合发展研究}
\author{石成玉}
\date{\today}

\begin{document}

\maketitle

\articleTitle

\begin{abstract}
	在“十五五”规划背景下，农业机械化作为农业现代化的关键支撑，正经历从“点线”成果迈向系统优势构建的转型。本文基于新质生产力理论，系统分析了“十五五”时期中国面临的战略机遇，梳理了农业机械化发展现状与阻碍因素，探讨了数字技术赋能农业产业链的路径。研究表明，农业机械化与农业现代化的深度融合，需要以新质生产力为引领，通过技术创新、要素重组、制度优化，实现农业全产业链的智能化、绿色化、可持续化发展。针对当前丘陵山区机械化滞后、数字技术应用不足、人才结构失衡等阻碍，提出强化科技攻关、深化数字赋能、优化制度环境等建议，为“十五五”时期推进农业强国建设提供理论支撑和实践参考。本文强调，农业机械化不仅是技术问题，更是实现中国式现代化的战略支点，其发展水平直接决定农业农村现代化的成色。
\end{abstract}

关键词：十五五规划；农业机械化；农业现代化；新质生产力；数字技术；农业服务贸易

\section{十五五计划}

\subsection{十五五时期中国的战略机遇}

“十五五”规划（2026-2030年）是承前启后、继往开来的关键五年，标志着中国向2035年基本实现社会主义现代化迈进的战略阶段。在这一时期，中国面临多重战略机遇，为农业机械化与农业现代化融合发展提供了历史性窗口。

首先，全球新一轮科技革命和产业变革为农业数字化转型提供强大引擎。中国机械工业联合会指出，“人工智能、大数据、物联网等数字技术与实体经济深度融合，为农业装备智能化升级创造了前所未有的机遇”\cite{BenKanBianJiBuMaoDingShiWuWuAILingXinCheng2025YangGuoQiShuZhiHuaZhuanXingYuZhiNengZhiZaoLunTanZaiJingZhaoKai2025}。数字技术正从“工具”向“核心要素”转变，推动农业从“经验驱动”转向“数据驱动”，为农机装备智能化提供技术基础。

其次，全球农业格局重塑带来农业服务贸易新机遇。胡冰川强调，“农业服务贸易将成为全球农业发展的重要增长极，将通过服务贸易的拓展，促进农业产业链的国际化布局，提升中国农业的国际竞争力”\cite{HuBingChuanNongYeFuWuMaoYiYuQuanQiuNongYeGeJuChongSuShiWuWuShiQiZhongGuoDeZhanLueJiYu2025}。“十五五”时期，中国可依托“一带一路”倡议，将农机装备出口、农业技术服务等纳入全球农业合作框架，实现从农产品贸易向农业服务贸易的升级。

第三，新质生产力发展为农业现代化注入新动能。马旭指出，“新质生产力是由农业科学技术革命性突破、各类农业生产要素创新性配置、农业产业深度转型升级催生而成的，以全要素生产率大幅提升为核心标志的优质生产力形态”\cite{MaXuXinZhiShengChanLiQuDongNongYeNongCunXianDaiHuaJiLiZuZhiYuYouHua2025}。这为农业机械化从“机械替代人力”向“智能优化生产”跃升提供了理论指引，使农业机械化成为发展新质生产力的重要载体。

\subsection{农业现代化}

“十五五”规划将农业现代化定位为“中国式现代化全局和成色”的关键支撑。叶兴庆系统阐述了“十五五”时期加快农业现代化的总体思路，指出新时代农业现代化是“用现代工业力量装备、用现代科学技术武装、以现代治理理论和方法经营、生产效率达到现代先进水平的农业”\cite{YeXingQingShenRuLingHuiShiWuWuShiQiJiaKuaiNongYeXianDaiHuaDeZongTiSiLu2025}。其核心表征包括：生产程序机械化、生产技术科学化、农业产业化、农业信息化、农业发展可持续化、劳动主体智能化。

农业现代化的内涵在“十五五”时期进一步深化。张洪霞强调，“农业新质生产力赋能农业现代化，要求突破传统机械化局限，实现从'机器替代人'到'智能赋能人'的转变”\cite{ZhangHongXiaNongYeXinZhiShengChanLiFuNengNongYeXianDaiHuaNeiHanYiYunXianShiKunJingJiShiJianLuJing2025}。这不仅体现在农机装备的智能化升级，更体现在农业生产全链条的数字化重构。例如，广西来宾市“滴滴农机”平台通过数字技术实现农机与农活的精准匹配，使作业效率提升30\%，交易成本降低25\%\cite{MaoMeiYanShuZiJiShuDuiNongYeChanYeLianXianDaiHuaDeKongJianYiChuXiaoYingJiYingXiangJiZhi2025}，这正是农业现代化从“单环节机械化”迈向“全链条智能化”的生动实践。

\section{农业机械的发展情况}

“十五五”规划明确要求“把农业建成现代化大产业”，农业机械化作为核心支撑，发展态势显著向好。

\subsection{机械化水平持续提升}

全国主要农作物耕种收综合机械化率已从“十三五”末的71\%提升至2025年的82\%，水稻、小麦、玉米等主要粮食作物耕种收综合机械化率分别达92\%、95\%、85\%\cite{FuHuaPingShenGengWoYeJiXianXingWoShiYiNongJiHuaGaiGeHuiJiuNongYeZhenXingXinTuJing2025}。保山市作为典型案例，通过农机化改革，建成3个万亩级全程机械化示范基地，农机总动力达45万千瓦，拥有各类机具12万台套，主要农作物耕种收综合机械化率达85.3\%\cite{FuHuaPingShenGengWoYeJiXianXingWoShiYiNongJiHuaGaiGeHuiJiuNongYeZhenXingXinTuJing2025}。建宁县数据显示，2024年农业机械总动力达19.64万千瓦，拥有各类机具4.6万台套，水稻耕种收综合机械化率达85.18\%\cite{FuHuaPingShenGengWoYeJiXianXingWoShiYiNongJiHuaGaiGeHuiJiuNongYeZhenXingXinTuJing2025}，充分验证了机械化水平的快速提升。

\subsection{丘陵山区农机装备突破}

针对丘陵山区机械化滞后问题，智能农机研发取得突破性进展。林凯报道的“智能农机'跑'上坡”案例显示，针对云南丘陵地区地形，研发的履带式自走式谷物收割机、山地轨道运输机等装备，使丘陵山区机械化率从不足40\%提升至65\%\cite{LinKaiZhiNengNongJiPaoShangPoQiuLingNongYeLuGengKuan2025}。这些装备通过北斗导航、智能感知技术，解决了坡度大、地块小的作业难题，为“十五五”时期丘陵山区机械化全覆盖奠定基础。

\subsection{数字技术深度融合}

数字技术正深度赋能农业机械化。中国机械工业联合会指出，“AI技术在农机装备中的应用已从辅助决策迈向自主作业，如智能拖拉机实现自动驾驶、变量施肥、精准喷药，作业效率提升30\%，资源利用率提高25\%”\cite{BenKanBianJiBuMaoDingShiWuWuAILingXinCheng2025YangGuoQiShuZhiHuaZhuanXingYuZhiNengZhiZaoLunTanZaiJingZhaoKai2025}。广西“滴滴农机”平台整合农机6000余台，注册用户超1000人，线上作业面积超1.2万亩，交易金额达430万元\cite{MaoMeiYanShuZiJiShuDuiNongYeChanYeLianXianDaiHuaDeKongJianYiChuXiaoYingJiYingXiangJiZhi2025}，标志着农业机械化从“单一装备应用”向“数字化服务生态”转型。

\subsection{农业机械化的阻碍}

尽管发展迅速，农业机械化仍面临多重阻碍，制约“十五五”规划目标实现。

\subsection{技术瓶颈与装备短板}

丘陵山区农机装备研发滞后是主要瓶颈。李金锴等指出，“我国丘陵山区农机装备自给率不足40\%，高端智能农机核心部件依赖进口，如导航系统、智能传感器等进口占比达65\%”\cite{LiJinKaiRenGongZhiNengFuNengNongYeLuSeZhuanXingDeZuoYongJiZhiXianShiKunJingYuTuiJinLuJing2025}。保山市案例显示，山区地块平均面积仅1.2亩，传统农机难以适应，导致机械化率比平原地区低25个百分点\cite{FuHuaPingShenGengWoYeJiXianXingWoShiYiNongJiHuaGaiGeHuiJiuNongYeZhenXingXinTuJing2025}。

\subsection{要素配置不均衡}

要素配置失衡制约机械化推广。马旭研究发现，“农业机械化发展面临'三缺'：缺资金（农机购置补贴仅覆盖30\%成本）、缺人才（农机操作员中大专以上文化占比不足15\%）、缺服务（农事服务中心覆盖不足40\%）”\cite{MaXuXinZhiShengChanLiQuDongNongYeNongCunXianDaiHuaJiLiZuZhiYuYouHua2025}。建宁县数据显示，农机服务组织中，专业人才占比仅18\%，导致技术应用效果打折扣。

\subsection{制度环境待优化}

制度环境制约机械化深度发展。魏凡森指出，“农机跨区作业协调机制不健全，跨省农机调度难；农机保险覆盖率不足30\%，制约农户购买意愿；农机与土地经营权衔接不畅，影响规模化应用”\cite{WeiFanSenNongYeYuQuYuJingJiRongHeFaZhanDeCuoShiTanTao2025}。“十五五”规划中，农业机械化与土地制度改革的协同不足，导致农机装备使用效率降低。

\subsection{农业服务贸易壁垒}

农业服务贸易国际壁垒影响农机出口。胡冰川强调，“欧盟、美国对农机技术标准设置高门槛，如欧盟农机排放标准比我国高30\%，导致我国农机出口受阻，2024年农机出口额仅占全球15\%”\cite{HuBingChuanNongYeFuWuMaoYiYuQuanQiuNongYeGeJuChongSuShiWuWuShiQiZhongGuoDeZhanLueJiYu2025}。这制约了农业机械化从国内走向国际的进程。

\section{促进农业机械化的建议}

针对上述阻碍，“十五五”时期应从技术、要素、制度、国际四个维度系统推进农业机械化发展。

\subsection{强化科技攻关与装备升级}

聚焦关键核心技术突破。建议设立“丘陵山区农机装备”国家科技重大专项，重点突破智能导航、精准作业、轻量化等技术，力争2028年前实现高端农机核心部件国产化率80\%以上\cite{LiJinKaiRenGongZhiNengFuNengNongYeLuSeZhuanXingDeZuoYongJiZhiXianShiKunJingYuTuiJinLuJing2025}。支持企业建立“产学研用”创新联合体，如保山市已建立“农机研发-示范应用-推广服务”闭环体系，使农机研发周期缩短40\%\cite{FuHuaPingShenGengWoYeJiXianXingWoShiYiNongJiHuaGaiGeHuiJiuNongYeZhenXingXinTuJing2025}。

\subsection{深化数字赋能与服务创新}

构建农业机械化数字生态。推广“滴滴农机”模式，建设全国农机调度平台，整合农机资源，实现“一键下单、智能匹配、全程追踪”。毛美燕指出，“数字技术对农业产业链现代化具有空间溢出效应，可使农机利用率提升25\%，服务半径扩大50\%”\cite{MaoMeiYanShuZiJiShuDuiNongYeChanYeLianXianDaiHuaDeKongJianYiChuXiaoYingJiYingXiangJiZhi2025}。同时，开发农机数字化培训平台，提升农民数字素养，目标“十五五”末农机操作员数字技能达标率达70\%。

\subsection{优化制度环境与要素配置}

完善农机发展制度体系。建议修订《农业机械化促进法》，建立“农机购置补贴+作业补贴+保险补贴”政策组合，将补贴覆盖比例提升至60\%\cite{MaXuXinZhiShengChanLiQuDongNongYeNongCunXianDaiHuaJiLiZuZhiYuYouHua2025}。推动土地经营权流转，建立“农机+土地”一体化运营机制，如建宁县推行“农机合作社+土地托管”模式，使农机使用效率提升35\%\cite{FuHuaPingShenGengWoYeJiXianXingWoShiYiNongJiHuaGaiGeHuiJiuNongYeZhenXingXinTuJing2025}。建立跨省农机作业协调机制，简化跨区作业手续。

\subsection{拓展农业服务贸易与国际合作}

提升农业服务贸易国际竞争力。支持农机企业“走出去”，参与国际标准制定，如推动中国农机标准纳入“一带一路”合作框架。胡冰川建议，“通过'农业服务贸易示范区'建设，培育5-10个国际知名农机服务品牌，将农机出口额提升至全球25\%”\cite{HuBingChuanNongYeFuWuMaoYiYuQuanQiuNongYeGeJuChongSuShiWuWuShiQiZhongGuoDeZhanLueJiYu2025}。同时，深化与东盟、非洲等区域合作，输出农机技术与服务，实现“中国农机”向“全球农机”的转变。
结语

“十五五”时期是农业机械化与农业现代化深度融合的关键窗口期。农业机械化不仅是技术问题，更是实现中国式现代化的战略支点，其发展水平直接决定农业农村现代化的成色。在新质生产力引领下，农业机械化正从“机械替代人力”向“智能赋能生产”跃升，从“单环节应用”向“全产业链协同”转型。通过强化科技攻关、深化数字赋能、优化制度环境、拓展国际合作，农业机械化将为“十五五”规划目标实现提供坚实支撑，为中国农业高质量发展注入新动能，为全球农业现代化贡献中国智慧和中国方案。

\printbibliography

\end{document}